\chapter{Active Galactic Nuclei: Accretion, Line Regions, and Outflows}
\label{chap:agn}

\section{The central engine as an energy-conversion system}

An active galactic nucleus (AGN) is not defined by one spectral line. It is a
coupled system in which gas loses angular momentum, releases gravitational
binding energy near a supermassive black hole, produces a broadband radiation
field, and sometimes launches winds and relativistic jets. The observables
depend on black-hole mass, accretion rate, spin, magnetic flux, viewing angle,
obscuration, time variability, and host-galaxy contamination. A taxonomy
based only on apparent line width or radio brightness therefore mixes
intrinsic physics with orientation and selection \citep{UrryPadovani1995}.

The Eddington luminosity follows from balancing electron-scattering radiation
force against gravity in ionized hydrogen:
\begin{equation}
 L_{\rm Edd}=\frac{4\pi GM_{\rm BH}m_{\rm p}c}{\sigma_{\rm T}}
 \simeq 1.26\times10^{38}
 \left(\frac{M_{\rm BH}}{M_\odot}\right)
 {\rm erg\,s^{-1}}.
 \label{eq:eddington}
\end{equation}
The Eddington ratio is $\lambda_{\rm Edd}=L_{\rm bol}/L_{\rm Edd}$. If the
radiative efficiency is $\eta$,
\begin{equation}
 L_{\rm bol}=\eta\dot{M}_{\rm acc}c^2.
 \label{eq:agn-accretion-luminosity}
\end{equation}
Neither $\eta$ nor the bolometric correction is universal. They depend on
accretion state, spin, spectral energy distribution (SED), and anisotropy.

In a geometrically thin, optically thick disk, the local effective temperature
approximately follows
\begin{equation}
 \sigma T_{\rm eff}^{4}(R)=
 \frac{3GM_{\rm BH}\dot{M}}{8\pi R^3}
 \left[1-\left(\frac{R_{\rm in}}{R}\right)^{1/2}\right],
 \label{eq:thin-disk}
\end{equation}
with relativistic corrections near the innermost stable orbit
\citep{ShakuraSunyaev1973}. Integrating blackbody-like annuli motivates the
optical/UV big blue bump. A hot, compact corona inverse-Compton scatters seed
photons into an X-ray power law. Dust outside its sublimation radius absorbs
UV/optical radiation and reradiates in the infrared. These components
illuminate different line-forming regions with an SED that is not directly
observed in every direction.

At very low Eddington ratio, radiatively inefficient flows can be hot and
geometrically thick; at high ratio, advection and radiation pressure modify
the thin-disk picture. Applying a thin-disk continuum fit outside its regime
can produce precise but model-specific masses and accretion rates.

\section{A scale hierarchy}

\begin{description}
 \item[Event-horizon and inner-disk scales:] gravitational radii
 $r_{\rm g}=GM_{\rm BH}/c^2$, probed by rapid variability, relativistic
 reflection, and sometimes interferometric imaging.
 \item[X-ray corona:] compact region producing the hard continuum; its
 geometry remains uncertain and may change with accretion state.
 \item[Broad-line region (BLR):] dense, rapidly moving gas producing permitted
 lines with widths of thousands of $\mathrm{km\,s^{-1}}$; forbidden lines are
 collisionally suppressed at high density.
 \item[Dusty structure:] an anisotropic, clumpy distribution near the dust
 sublimation scale, not necessarily a smooth torus.
 \item[Narrow-line region (NLR):] lower-density gas extending from tens of
 parsecs to kiloparsecs, producing forbidden and permitted lines with
 narrower profiles.
 \item[Radio jet and lobes:] magnetized relativistic flow spanning sub-parsec
 to megaparsec scales in the most powerful sources.
 \item[Galaxy-scale outflow:] ionized, neutral, atomic, and molecular gas
 accelerated or pressurized by radiation, winds, jets, star formation, or
 their combination.
\end{description}

This hierarchy matters because apertures mix it. A fiber spectrum may combine
BLR light, NLR emission, H II regions, diffuse ionized gas, shocks, and stellar
continuum. An X-ray spectrum can mix coronal continuum, disk reflection,
distant reflection, warm absorption, soft excess, and host emission.
Spectral decomposition is therefore a spatial model even when the object is
unresolved.

\section{Emission-line regions}

\subsection{Permitted broad lines}

Hydrogen Balmer lines, Mg II, C IV, Ly$\alpha$, and other permitted
transitions dominate many type-1 AGN spectra. Their profiles encode a
luminosity-weighted transfer function rather than a simple orbit. Rotation,
inflow, outflow, anisotropic emission, obscuration, responsivity, and
radiation pressure all contribute. C IV is particularly susceptible to a
blueshifted wind component; treating its full width as purely virial can bias
single-epoch black-hole masses.

The virial mass estimator is
\begin{equation}
 M_{\rm BH}=f_{\rm vir}\frac{R_{\rm BLR}\Delta V^2}{G},
 \label{eq:virial-mass}
\end{equation}
where $\Delta V$ is a line-width measure and $f_{\rm vir}$ absorbs BLR
geometry, inclination, and kinematics. Different use of FWHM versus line
dispersion and mean versus rms spectra changes the calibration. The result is
not a direct dynamical mass unless $R_{\rm BLR}$ and the transfer geometry are
constrained.

\subsection{Reverberation mapping}

If a variable continuum drives line variability, the observed line light
curve is
\begin{equation}
 L_{\rm line}(v,t)=
 \int \Psi(v,\tau)\,C(t-\tau)\,{\rm d}\tau,
 \label{eq:reverberation}
\end{equation}
where $\Psi(v,\tau)$ is the velocity-delay map
\citep{BlandfordMcKee1982,Peterson1993}. A characteristic lag gives
$R_{\rm BLR}\simeq c\tau$ in the rest frame. Velocity-resolved lags can
distinguish idealized rotation, inflow, and outflow patterns, although finite
cadence, seasonal gaps, red-noise variability, and line blending blur the
response.

Campaign design must match three timescales: sampling cadence shorter than the
expected lag structure, duration several times the largest lag, and enough
continuum variability to drive a measurable response. The observed lag is
dilated by $(1+z)$. Flux calibration must be stable across epochs; aperture
changes and seeing-dependent host contamination can imitate variability.

\subsection{Narrow-line diagnostics}

Classical diagnostic diagrams compare lines close in wavelength so that dust
and flux-calibration effects are modest. The most familiar plane is
\begin{equation}
 \log\frac{[\mathrm{O\,III}]\,5007}{\mathrm{H}\beta}
 \quad {\rm versus}\quad
 \log\frac{[\mathrm{N\,II}]\,6583}{\mathrm{H}\alpha},
 \label{eq:bpt}
\end{equation}
with companion planes using [S II]/H$\alpha$ and [O I]/H$\alpha$
\citep{Baldwin1981,Kewley2001,Kauffmann2003,Kewley2006}. Hard AGN continua
extend partially ionized zones and enhance low-ionization lines, while stellar
H II regions occupy a sequence controlled by metallicity and ionization
parameter.

The boundaries are not laws of nature. They depend on abundance pattern,
stellar populations, redshift, density, diffuse ionized gas, shocks, and
aperture mixing. At high redshift, star-forming galaxies can shift relative
to local training samples. A spatially resolved mixing sequence is often more
informative than one integrated classification.

Density-sensitive doublets such as [S II] 6716/6731 and [O II] 3726/3729
constrain electron density only over their responsive range and for the zone
emitting the ions. Temperature-sensitive auroral-to-nebular ratios such as
[O III] 4363/(4959+5007) measure a strongly emissivity-weighted temperature.
In an inhomogeneous NLR, a single density and temperature are summaries, not
a full structure \citep{OsterbrockFerland2006}.

\section{Obscuration and orientation}

The basic unification insight is that an intrinsically similar central engine
can look different when the BLR is directly visible, hidden by dusty material,
or seen through a relativistically beamed jet \citep{UrryPadovani1995}.
Polarized broad lines, ionization cones, radio morphology, and infrared
emission support this view. Modern unification is statistical rather than
absolute. Covering factor, luminosity, host dust, accretion state, and time
variability also matter; some changing-look AGN alter spectral type on human
timescales.

Line ratios involving widely separated wavelengths require an extinction
model. For a foreground-screen approximation,
\begin{equation}
 F_{\rm obs}(\lambda)=F_{\rm int}(\lambda)
 10^{-0.4E(B-V)k(\lambda)}.
 \label{eq:reddening}
\end{equation}
The Balmer decrement can constrain reddening, but intrinsic ratios vary with
density, optical depth, collisional excitation, and radiative transfer.
Mixed dust-emitter geometry does not obey a simple screen.

\section{Absorption from nuclear winds}

\subsection{UV absorption classes}

AGN spectra show narrow absorption lines, mini-broad troughs, and broad
absorption lines (BALs). Common ions include C IV, Si IV, N V, O VI, Mg II,
and metastable or excited transitions. Blueshift establishes motion toward
the observer relative to the systemic frame, but distance and mass flux remain
unknown without a density or geometry constraint.

Doublet residuals often reveal saturation plus partial coverage. Excited-state
to ground-state ratios can constrain density; combining density with the
ionization parameter gives a radial estimate
\begin{equation}
 r=\left(\frac{Q_{\rm H}}{4\pi c\,U n_{\rm H}}\right)^{1/2},
 \label{eq:agn-wind-radius}
\end{equation}
where $Q_{\rm H}$ is the ionizing photon rate. The result inherits uncertainty
from the unobserved extreme-UV SED, level populations, shielding, and whether
the transitions share a phase.

Variability supplies a complementary constraint. A response to continuum
changes can bound recombination time and density; transverse changes in
covering can instead reflect cloud motion. Distinguishing ionization change
from occultation requires multiple ions and a measured ionizing continuum.

\subsection{X-ray warm absorbers and ultra-fast outflows}

Soft-X-ray spectra of many AGN contain photoionized absorption from oxygen,
neon, iron, and other elements, often in multiple ionization and velocity
components. At Fe K energies, blueshifted Fe XXV and Fe XXVI absorption can
indicate velocities near a tenth of the speed of light in ultra-fast outflows
\citep{Tombesi2010,KingPounds2015}. The identification is demanding: the
continuum and reflection model, detector features, blind energy search, and
number of trials all affect significance. Confirmation through multiple
transitions or epochs is more persuasive than one residual.

X-ray ionization is commonly described by
\begin{equation}
 \xi=\frac{L_{\rm ion}}{n_{\rm H}r^2}
 \quad [{\rm erg\,cm\,s^{-1}}],
 \label{eq:xi}
\end{equation}
over a stated ionizing energy range. Conventions for $L_{\rm ion}$ differ, so
$\xi$ values from different codes or papers are not automatically
interchangeable. A distribution of column over ionization parameter may be
more realistic than a small number of discrete slabs.

\section{Outflow mass, momentum, and power}

For a thin partial shell of radius $R$, column $N_{\rm H}$, solid-angle
fraction $\Omega/4\pi$, and characteristic speed $v$, a widely used estimate
is
\begin{equation}
 \dot{M}_{\rm out}\simeq
 \Omega\,R\,\mu m_{\rm p}N_{\rm H}v,
 \label{eq:outflow-rate}
\end{equation}
where here $\Omega$ is in steradians. Authors sometimes absorb $4\pi$ into a
dimensionless covering factor; the convention must be stated. The associated
momentum and kinetic-power rates are
\begin{equation}
 \dot{p}_{\rm out}=\dot{M}_{\rm out}v,\qquad
 \dot{E}_{\rm kin}=\frac{1}{2}\dot{M}_{\rm out}v^2.
 \label{eq:outflow-power}
\end{equation}

Another common estimator uses $M_{\rm out}v/R$ or $3M_{\rm out}v/R$ depending
on whether the assumed flow is a shell or a volume-filling, steady cone.
This factor is not bookkeeping trivia: it encodes a different history and
density profile. Radius, deprojection, electron density, conversion factor,
and low-velocity separation from the host can each dominate uncertainty.
Tables comparing outflow energetics should homogenize these choices before
claiming physical trends \citep{Veilleux2005,Fiore2017}.

Ionized mass from a recombination line scales roughly as
\begin{equation}
 M_{\rm ion}\propto
 \frac{L_{\rm recomb}}{n_{\rm e}\alpha_{\rm eff}h\nu},
\end{equation}
so adopting an electron density that is too low inflates the mass. Density
doublets may preferentially weight bright dense clumps, while most mass could
occupy a lower-density component. Molecular mass from CO requires an
$\alpha_{\rm CO}$ conversion; neutral atomic mass from absorption requires a
spin temperature and covering factor. Multi-phase outflow budgets are
therefore sums of differently conditioned estimates, not directly measured
totals \citep{Bolatto2013,Cicone2014}.

\section{Jets and radio diagnostics}

Synchrotron emission identifies relativistic electrons and magnetic fields.
For $S_\nu\propto\nu^{\alpha}$, a steep spectrum generally characterizes
optically thin lobes while a flat compact core can be a superposition of
self-absorbed jet components. Doppler beaming changes observed flux,
timescales, and apparent speed. Polarization and Faraday rotation probe
ordered fields and intervening magnetized plasma, but depolarization can be
internal or external.

Jet mechanical power is often inferred from cavities in hot X-ray gas:
\begin{equation}
 P_{\rm cav}\simeq \frac{E_{\rm cav}}{t_{\rm age}},
 \qquad
 E_{\rm cav}\simeq \frac{\gamma}{\gamma-1}pV.
\end{equation}
For relativistic plasma, the enthalpy is commonly approximated as $4pV$.
Buoyancy, sound-crossing, and refill times give different ages. Radio
luminosity is an imperfect proxy because synchrotron efficiency, spectral age,
and environment vary.

Jets can drive shocks, entrain cool gas, inflate bubbles, or maintain
turbulence. A spatial coincidence between radio and line emission is
suggestive, but velocity-resolved line ratios and pressure estimates are
needed to distinguish shock excitation, compression-enhanced star formation,
and AGN photoionization. The Blandford--Znajek and Blandford--Payne mechanisms
provide limiting physical pictures for spin-powered and disk-launched
outflows \citep{BlandfordZnajek1977,BlandfordPayne1982}.

\section{X-ray reflection and the Fe K complex}

Hard coronal photons striking optically thick material produce a reflected
continuum, absorption edges, and fluorescent lines. Near-neutral iron emits
Fe K$\alpha$ near 6.4 keV; He-like and H-like iron produce features near
6.7 and 6.97 keV. Reflection from the inner disk is shifted and skewed by
Doppler motion, gravitational redshift, and light bending
\citep{GeorgeFabian1991,Reynolds2014}.

Spin inference uses the inner radius implied by the full relativistic
reflection spectrum, not merely the red wing of one line. The fit depends on
disk ionization, density, elemental abundance, inclination, coronal
illumination profile, continuum curvature, absorption, and the assumption
that the disk reaches the innermost stable circular orbit. Robust analyses
compare alternative continua and reflection codes, fit a broad energy band,
and test whether spin is constrained by the data rather than a parameter
boundary.

A narrow 6.4-keV component can arise in distant Compton-thick material.
The Compton shoulder, K$\beta$ line, reflection hump, and variability help
locate it. In Compton-thick AGN, the direct continuum can be strongly
suppressed, making equivalent widths large and naive bolometric corrections
unreliable.

\section{Feedback across phases and timescales}

Nuclear winds can begin at sub-parsec scales while molecular outflows are
measured over hundreds of parsecs or kiloparsecs. Equality of their
instantaneous energy fluxes is not required: the phases are observed at
different epochs of an intermittent source. In an ideal momentum-driven
picture, shocked wind energy cools efficiently and the large-scale momentum
rate is of order $L/c$. In an ideal energy-conserving bubble, retained thermal
energy can produce a larger momentum flux \citep{KingPounds2015}. Leakage,
gravity, radiative losses, clumpy coupling, and source variability blur this
distinction.

Evidence for feedback should connect at least three levels:
\begin{enumerate}
  \item a credible AGN energy or momentum source;
  \item a kinematically and spatially resolved response in surrounding gas;
  \item an effect on cooling, star formation, or gas retention beyond what
  stellar feedback and gravitational motions plausibly explain.
\end{enumerate}
A broad line alone demonstrates disturbed gas, not quenching. Conversely,
maintenance-mode heating may prevent cooling without an obvious fast outflow.

\section{An AGN line-analysis checklist}

\begin{enumerate}
 \item Establish systemic redshift with a documented tracer.
 \item Separate instrumental resolution from intrinsic line width.
 \item Fit stellar continuum and Balmer absorption before narrow-line ratios.
 \item Decompose narrow and broad components jointly across linked lines,
 respecting fixed atomic ratios only where justified.
 \item Test reddening, aperture mixing, and spatially resolved diagnostic
 ratios.
 \item For absorption, fit doublets with optical depth and covering fraction;
 search all plausible blends.
 \item For mass outflow rates, publish the geometry, radius, density or
 conversion factor, projection convention, and velocity statistic.
 \item For X-ray residuals, include continuum and reflection alternatives and
 the blind-search trials factor.
 \item Compare phases at matched aperture and timescale where possible.
 \item Distinguish measured luminosity or column from inferred mass, power,
 and causal feedback.
\end{enumerate}

AGN spectra are richest when treated as a connected system. The UV line that
measures a nuclear wind, the optical lines that map an ionization cone, the
CO wings that reveal accelerated molecular gas, the radio lobes that carry
mechanical energy, and the X-ray lines that probe the corona and fastest wind
are not competing descriptions. They are different density, temperature,
radius, and time weightings of one evolving engine.
